\input{../../.lake/build/blueprint/library/Goldbach.tex}

\chapter{Two routes toward Goldbach}

This project formalizes Chen's $1+2$ theorem and Li--Liu's $1+1.9$ theorem,
while building reusable analytic number theory toward stronger Goldbach results.
The two developments retain independent public interfaces and proof routes.
Use \texttt{Goldbach} for Chen, \texttt{Goldbach.OnePlusOneNine} for Li--Liu,
and \texttt{Goldbach.All} to import both.

Chen's theorem represents every sufficiently large even natural number as
\[
N=p+q,
\]
where $p$ is prime and $q$ is either prime or a product of two primes.
The two factors in the product may coincide. Li--Liu's theorem also controls
the relative sizes of the factors:
\[
N=p+rq,\qquad p,q\text{ prime},\qquad r=1\text{ or }r\text{ prime},
\qquad r^{10}\le q^9.
\]
The last condition is equivalently $r\le q^{9/10}$.
Both existence statements have an existential threshold.

\section{How to read the graph}

The interactive graph follows sixteen selected declarations: three reusable
foundations, the seven-node Chen endpoint assembly, and six nodes in the
Li--Liu route. LeanArchitect extracts dependencies from compiled declarations;
arrows run from a dependency to its consumer. Each node has a mathematical
summary and a link to its exact Lean source. The selected graph is a reading
map through the full proof development. Open the Lean documentation to inspect
individual types, module imports and the surrounding lemmas.

The Li--Liu graph distinguishes the positive-margin route used by the public
existence theorem from the sharper weighted-count estimate used by the
quantitative theorem. This distinction follows the actual implementation.

\chapter{Reusable analytic foundations}

Several recurring arguments now have a common proof, with concrete applications
supplying their own arithmetic or geometric data. These interfaces connect the
existing developments and provide building blocks for further work.

\section{Transporting a finite prime-pair sum}

Liu's source separates the two prime-factor ranges. An admissible product
therefore determines its ordered prime pair uniquely. This permits a single
change-of-index theorem for arbitrary cutoffs, filters and additive kernels.
The main-term sum, the noncoprime correction, the arbitrary-kernel identity
and the pair-cardinality calculation all use this theorem. The filter and
cutoff remain part of each application's exact finite carrier.
\inputleannode{thm:pair-transport}

\section{Absorbing a logarithmic remainder}

Fix a positive logarithmic saving and a positive common lower bound for the
normalizing weight. The extra logarithmic power makes the error coefficient
tend to zero. After choosing one threshold, the estimate holds for every
normalization exponent and every weight above that lower bound.
Chen's singular-series estimates and Li--Liu's normalized terms use this
common argument. Each application retains its own error budget and cutoff.
\inputleannode{thm:log-saving}

\section{Paying for a grid boundary}

On the intersection of the grid region and the source region, a local bound
controls the integrand error. Outside the source region, finitely many boundary
strips cover the excess. Integrating these bounds yields one estimate that
allows the strips to overlap. The B8, B9, B10 and Liu prime-pair applications
supply their own strips, measures and kernel bounds. Their logarithmic
rectangle masses and endpoint conventions are retained in those applications.
\inputleannode{thm:grid-error}

\chapter{Chen's prime-plus-almost-prime theorem}

The Chen assembly combines a weighted lower sieve with the proved Liu--Pan
switched-source distribution estimate. The intermediate theorems expose that
distribution estimate as an input. Supplying its proved instance closes the
existence and quantitative conclusions.

\section{Distribution and assembly}
\inputleannode{thm:liu-pan}
\inputleannode{thm:conditional-count}
\inputleannode{thm:conditional-chen}

\section{Unconditional endpoints}
\inputleannode{thm:count-internal}
\inputleannode{thm:chen-internal}

\section{Public results}
\inputleannode{thm:representation}
\inputleannode{thm:chen}

\chapter{Li--Liu's constrained-factor theorem}

The finite weights detect primes that admit the required constrained-factor
representation. Sieve estimates then turn those finite inequalities into
positive lower bounds for sufficiently large even integers.
\inputleannode{thm:liliu-finite}

\section{Existence from a positive margin}

Certified bounds for the sieve integrals establish a positive margin. The
finite detection theorem converts that positivity into an eligible prime and
hence the public $N=p+rq$ representation.
\inputleannode{thm:liliu-margin}
\inputleannode{thm:liliu-existence}

\section{A quantitative count of distinct primes}

Write $D_{1,19/10}(N)$ for the number of eligible primes $p$, counting each
prime once even when it has several witnessing factorizations. The
normalization is the Liu singular series
\[
\mathfrak S_{\mathrm{Liu}}(N)=
\prod_{\substack{\ell>2\ \mathrm{prime}\\\ell\mid N}}
\frac{\ell-1}{\ell-2}
\prod_{\ell>2\ \mathrm{prime}}\left(1-\frac{1}{(\ell-1)^2}\right).
\]
The quantitative route applies the author G11 estimate and the certified
remaining estimates to their original counts before combining the signed
terms. It supplies, for every fixed real $\kappa<515093/800000000$, a threshold
beyond which
\[
\kappa\,\mathfrak S_{\mathrm{Liu}}(N)\frac{N}{(\log N)^2}
\le D_{1,19/10}(N)
\]
for all even $N$. The coefficient is fixed before the threshold is chosen.
\inputleannode{thm:liliu-ledger}
\inputleannode{thm:liliu-family}

Choosing a certified coefficient above $1/2500$ and using positivity of the
normalizing scale gives the paper's strict bound
\[
\frac{1}{2500}\,\mathfrak S_{\mathrm{Liu}}(N)\frac{N}{(\log N)^2}
< D_{1,19/10}(N),\qquad 1/2500=0.0004.
\]
\inputleannode{thm:liliu-count}

\chapter{Sources and verification}

The development builds on UyNewNas/chen-theorem-lean and the attributed
analytic number theory sources, with upstream copyright notices retained
under Apache-2.0. The provenance guide records these sources and the
Li--Liu mathematical inputs.

The repository contains the literal specifications, full proofs, theorem
guide and verification scripts. Both completed routes use Lean's standard
logical foundations: propositional extensionality, classical choice and
quotient soundness. Verification combines strict compilation, exact statement
and axiom checks, and independent kernel replay. The graph supplies a route
from the mathematical outline to the declarations checked by Lean.

Existing checkouts can retain \texttt{.lake/} when updating. With the same
package, toolchain and dependency pins, Lake reuses unchanged artifacts and
rebuilds affected modules. The README provides focused commands for each
public theorem and for the combined interface.

\url{https://github.com/subfish-zhou/goldbach-lean}
